\documentclass[11pt]{article}

% ─── Packages ────────────────────────────────────────────────────────────────
\usepackage{amsmath,amssymb,amsthm}
\usepackage[margin=1in]{geometry}
\usepackage{hyperref}
\usepackage{graphicx}
\usepackage{booktabs}
\usepackage{array}
\usepackage{multirow}
\usepackage{xcolor}
\usepackage{authblk}
\usepackage{setspace}
\usepackage{caption}
\usepackage{subcaption}
\usepackage{enumitem}
\usepackage{microtype}
\usepackage{siunitx}

\hypersetup{
  colorlinks=true,
  linkcolor=blue,
  citecolor=blue,
  urlcolor=blue
}

% ─── Theorem environments ────────────────────────────────────────────────────
\theoremstyle{plain}
\newtheorem{theorem}{Theorem}[section]
\newtheorem{lemma}[theorem]{Lemma}
\newtheorem{proposition}[theorem]{Proposition}
\newtheorem{corollary}[theorem]{Corollary}
\theoremstyle{definition}
\newtheorem{definition}{Definition}[section]
\theoremstyle{remark}
\newtheorem{remark}{Remark}[section]

% ─── Title block ─────────────────────────────────────────────────────────────
\title{\textbf{The Arctic Optimum: Temperature-Responsive\\
Coordination Entropy and the Design of Electrolytes\\
That Improve as Temperature Falls}}

\author[1]{Eric Robert Lawson}

\affil[1]{Independent Researcher, OrganismCore Initiative\\
\texttt{ORCID: 0009-0002-0414-6544}\\
Repository: \url{https://github.com/Eric-Robert-Lawson/attractor-oncology}}

\date{April 4, 2026\\
\small Pre-registration record:
\url{https://github.com/Eric-Robert-Lawson/attractor-oncology}\\
\small (timestamped commit history constitutes the priority
and pre-registration record)\\
\small Part of the SCE Framework preprint series.
Preprint~1 (foundation and fixed point):
same repository.\\
\small The full geometric structure of this series,
including the relationship between all candidate formulations,
mechanisms, and derivations, is navigable from the repository
index.}

% ─── Document ────────────────────────────────────────────────────────────────
\begin{document}

\maketitle

% ─── Abstract ────────────────────────────────────────────────────────────────
\begin{abstract}
The SCE Framework establishes $\mathrm{SCE}^{*} = 1.466$ as the
Shannon Coordination Entropy of the Li$^+$ first solvation shell at
which combined room-temperature (RT) and low-temperature (LT)
Coulombic efficiency (CE) is simultaneously maximised.  Prior papers
in this series derived three mechanisms for reaching
$\mathrm{SCE}^{*}$ at a fixed temperature: ESP matching, anion
reception, and within-class geometric diversity.  All three treat
$\mathrm{SCE}^{*}$ as a static target at 25\,$^\circ$C.

The present paper introduces a fundamentally different design
principle: the \textbf{temperature-responsive SCE gap}.  Rather than
designing an electrolyte to sit at $\mathrm{SCE}^{*}$ at
25\,$^\circ$C and then measuring whether LT performance follows,
we design an electrolyte whose SCE is a function of temperature
--- specifically, one whose SCE moves \emph{toward} $\mathrm{SCE}^{*}$
as temperature decreases.  At the operating temperature where the SCE
gap closes, LT performance is maximised.  The design target is not a
fixed-point value but a temperature-SCE trajectory that crosses
$\mathrm{SCE}^{*}$ at the intended operating temperature.

This principle defines the \textbf{Arctic Optimum}: a formulation
regime in which decreasing temperature drives the solvation structure
toward the fixed point rather than away from it.  The canonical
realisation is LiFSI 1.2~M in DME/TMP/TMB, a three-component system
combining the ESP-matching mechanism (DME/TMP, Preprint~2) with the
anion receptor mechanism (TMB, Preprint~3) in a single formulation
whose temperature coefficient of SCE is engineered to be negative:
SCE decreases as temperature decreases, approaching $\mathrm{SCE}^{*}$
from above at low temperature.

Arctic Candidate~A (LiFSI 1.2~M in DME/TMP/TMB, 3:1:0.5 v/v/v,
operating target $-30\,^\circ$C) is derived with exact formulation
rationale, temperature-SCE trajectory prediction,
performance predictions at RT and at $-30\,^\circ$C, and full
pre-registered falsification conditions.  The SCE ceiling
(Derivation~5 of the pre-registered record) is derived and confirmed:
above a maximum achievable SCE, further entropy increase is
structurally inaccessible in ether-based systems.  The arctic optimum
lies below this ceiling and above $\mathrm{SCE}^{*}$, defining the
accessible design space for temperature-responsive electrolytes.

The derivation record, MU session confirmation, and full falsification
conditions are archived at
\url{https://github.com/Eric-Robert-Lawson/attractor-oncology}.
\end{abstract}

\vspace{1em}
\noindent\textbf{Keywords:} lithium metal battery, electrolyte
design, arctic, low temperature, temperature-responsive,
coordination entropy, SCE*, DME, TMP, TMB, solvation structure,
fixed point, SCE ceiling

\tableofcontents
\newpage

% ─── 1. Introduction ─────────────────────────────────────────────────────────
\section{Introduction}

\subsection{The Fixed Point and the Temperature Problem}

The SCE Framework (Preprint~1) establishes that:

\begin{equation}
  \mathrm{SCE}^{*} = 1.466
  \label{eq:fixed_point}
\end{equation}

maximises combined RT and LT Coulombic efficiency simultaneously.
The framework equations are:

\begin{align}
  \mathrm{CE}_{\mathrm{RT}} &= 100.1 - e^{1.493 + 1.230 \cdot \mathrm{SCE}}
  \label{eq:CE_RT}\\[4pt]
  \mathrm{CE}_{\mathrm{LT}} &= 11.91 + 33.21 \cdot \mathrm{SCE}
  \label{eq:CE_LT}
\end{align}

with $R^2 = 0.828$, $p = 4.5 \times 10^{-6}$, confirmed across 21
systems.

A critical implicit assumption underlies all prior papers in this
series: that SCE is measured and optimised at a fixed reference
temperature (25\,$^\circ$C) and that the resulting formulation then
performs at low temperature.  This assumption is valid for
applications where the operating temperature is close to 25\,$^\circ$C.
It is not optimal for applications where the intended operating
temperature is substantially lower --- specifically, in the range
$-20\,^\circ$C to $-40\,^\circ$C characteristic of arctic environments,
cold storage systems, and high-altitude aerospace applications.

\subsection{The Temperature-SCE Coupling}

Solvation structure is temperature-dependent.  As temperature
decreases:
\begin{itemize}[leftmargin=2em]
  \item Thermal energy $k_BT$ decreases.
  \item The Boltzmann weighting of higher-energy coordination
        configurations decreases.
  \item The population distribution narrows: $p_{\max}$ rises and
        minority configurations are depopulated.
  \item SCE decreases in most systems.
\end{itemize}

This is the standard direction: cooling orders the solvation shell.
Most electrolyte systems approach lower SCE as temperature falls ---
moving \emph{away} from $\mathrm{SCE}^{*}$ if they were at or above
the fixed point, or deeper below it if they were already below.

But the direction of the temperature-SCE coupling is not fixed by
physics.  It is determined by the specific balance of competing
coordination mechanisms in a given formulation.  In a system where
the dominant coordination mode at RT involves two competing geometries
of similar energy, cooling preferentially suppresses the higher-energy
geometry, concentrating the population into the lower-energy mode and
\emph{reducing} SCE.  But in a system where two coordination modes
compete and one is entropically stabilised at RT while the other is
enthalpically stabilised at low temperature, cooling can shift the
population \emph{toward} the enthalpically preferred mode, changing
the distribution in a direction that may raise or maintain SCE
depending on the geometry.

The key insight is this:

\begin{definition}[Temperature-responsive SCE gap]
A \textbf{temperature-responsive SCE gap} exists in a formulation
when the SCE is a monotonically decreasing (or non-monotonic)
function of temperature such that there exists a target temperature
$T^{*} < 25\,^\circ\mathrm{C}$ at which:
\begin{equation}
  \mathrm{SCE}(T^{*}) = \mathrm{SCE}^{*} = 1.466
  \label{eq:T_target}
\end{equation}
Above $T^{*}$, $\mathrm{SCE} > \mathrm{SCE}^{*}$ (shell too
disordered at RT).  As temperature falls toward $T^{*}$, SCE
decreases toward the fixed point.  At $T^{*}$, the fixed point
is reached.  Below $T^{*}$, SCE falls below $\mathrm{SCE}^{*}$
(shell too ordered at the lowest temperatures).
\end{definition}

A formulation exhibiting a temperature-responsive SCE gap does not
need to be at $\mathrm{SCE}^{*}$ at 25\,$^\circ$C.  It needs to
\emph{cross} $\mathrm{SCE}^{*}$ at the intended operating temperature.
This is a fundamentally different design paradigm.

\subsection{The Arctic Optimum}

The \textbf{arctic optimum} is the formulation regime in which the
temperature-responsive SCE gap closes at a target operating
temperature $T^{*}$ in the range $-20\,^\circ$C to $-40\,^\circ$C.
An electrolyte designed for the arctic optimum:
\begin{itemize}[leftmargin=2em]
  \item Has $\mathrm{SCE}_\mathrm{RT} > \mathrm{SCE}^{*}$: at
        25\,$^\circ$C the shell is slightly too disordered.
  \item Has a negative temperature coefficient of SCE:
        $\mathrm{d}(\mathrm{SCE})/\mathrm{d}T > 0$
        (SCE decreases as temperature decreases).
  \item Crosses $\mathrm{SCE}^{*}$ at $T^{*}$: at the target
        operating temperature the shell reaches the fixed point.
  \item Has $\mathrm{SCE}_{T^*} = \mathrm{SCE}^{*}$: at the
        intended operating temperature, peak dual performance is
        achieved.
\end{itemize}

This is the design principle of Arctic Candidate~A.

\subsection{Why the Three-Component System}

The three mechanisms established in Preprints~2, 3, and~4 operate
at fixed temperature.  Arctic Candidate~A combines two of them in a
single formulation:

\begin{enumerate}[leftmargin=2em]
  \item \textbf{ESP matching} (DME/TMP): raises SCE from below
        $\mathrm{SCE}^{*}$ at RT.  The cross-class mechanism
        produces a population split whose temperature coefficient
        depends on the relative binding enthalpies of TMP and DME
        to Li$^+$.
  \item \textbf{Anion reception} (TMB): provides a temperature-
        sensitive anion sequestration mechanism.  At low temperature,
        the B--N Lewis acid--base interaction strengthens (entropy
        of activation becomes less unfavourable), sequestering more
        FSI$^-$ and reducing the CIP/AGG contribution to SCE.
\end{enumerate}

The combination of these two mechanisms in a single formulation
produces a system whose overall SCE temperature coefficient can be
tuned to place the SCE--temperature crossing at a target operating
temperature.

\subsection{Structure of This Paper}

Section~\ref{sec:SCE_ceiling} derives the SCE ceiling: the maximum
achievable SCE in ether-based Li metal electrolyte systems.
Section~\ref{sec:arctic_design_space} establishes the arctic design
space between the ceiling and $\mathrm{SCE}^{*}$.
Section~\ref{sec:temperature_coefficient} derives the
temperature coefficient of SCE for the three-component system.
Section~\ref{sec:candidateA} states Arctic Candidate~A with the
full formulation rationale, trajectory prediction, and performance
predictions.  Section~\ref{sec:falsification} states the
pre-registered falsification conditions.
Section~\ref{sec:literature} reports the arctic literature
confirmation.  Section~\ref{sec:discussion} discusses implications
and limitations.

% ─── 2. The SCE Ceiling ───────────────────────────────────────────────────────
\section{The SCE Ceiling}
\label{sec:SCE_ceiling}

\subsection{Derivation}

The Shannon Coordination Entropy of the Li$^+$ first solvation shell
is bounded above by the entropy of the maximally disordered
configuration space accessible to the system:

\begin{equation}
  \mathrm{SCE}_{\max} = \ln(n_{\max})
  \label{eq:SCE_max_def}
\end{equation}

where $n_{\max}$ is the number of distinct, meaningfully occupied
coordination configurations accessible to Li$^+$ in the given
solvent system.

For ether-based Li metal electrolytes at moderate salt concentrations
(0.5--2.0~M LiFSI), the coordination configurations are bounded by:
\begin{itemize}[leftmargin=2em]
  \item SSIP configurations: classified by solvent identity (DME,
        TMP, 2-MeTHF, DOL, etc.) and coordination mode (bidentate,
        monodentate, chelate).  Maximum $n_\mathrm{SSIP} \approx 4$
        distinct configurations for a three-component solvent system.
  \item CIP configurations: Li$^+$ with one FSI$^-$ in shell.
        Maximum $n_\mathrm{CIP} \approx 2$ (two FSI$^-$ orientations).
  \item AGG configurations: Li$^+$ with two FSI$^-$ in shell.
        Maximum $n_\mathrm{AGG} \approx 1$ at moderate
        concentrations.
\end{itemize}

Maximum accessible $n_{\max} \approx 7$ configurations in a
three-component ether system at moderate LiFSI concentration.

\begin{equation}
  \mathrm{SCE}_{\mathrm{ceiling}} = \ln(7) \approx 1.946
  \label{eq:SCE_ceiling}
\end{equation}

However, this ceiling is not achievable: it requires all seven
configurations to be equally populated, which is never the case
in a physically realisable system with non-zero enthalpic
preferences between coordination modes.  A practical ceiling
accounting for realistic population distributions:

\begin{equation}
  \mathrm{SCE}_{\mathrm{practical~ceiling}} \approx 1.75
  \label{eq:SCE_practical_ceiling}
\end{equation}

This practical ceiling is derived in Derivation~5 of the
pre-registered record by computing the maximum SCE achievable when
the dominant configuration fraction is constrained to
$p_{\max} \geq 0.20$ (a physically realistic minimum for any
system with enthalpic preferences).

\subsection{The Design Space}

The arctic design space is defined by:

\begin{equation}
  \mathrm{SCE}^{*} < \mathrm{SCE}_\mathrm{RT} \leq
  \mathrm{SCE}_\mathrm{practical~ceiling}
  \quad \Rightarrow \quad
  1.466 < \mathrm{SCE}_\mathrm{RT} \leq 1.75
  \label{eq:arctic_design_space}
\end{equation}

A formulation with RT SCE in this range and a negative temperature
coefficient of SCE will cross $\mathrm{SCE}^{*}$ at some temperature
below 25\,$^\circ$C.  The target temperature $T^{*}$ is set by
selecting the RT SCE value and the magnitude of the temperature
coefficient:

\begin{equation}
  T^{*} = 25 - \frac{\mathrm{SCE}_\mathrm{RT} - \mathrm{SCE}^{*}}
  {|\mathrm{d}(\mathrm{SCE})/\mathrm{d}T|}
  \label{eq:T_star_formula}
\end{equation}

where the temperature coefficient
$|\mathrm{d}(\mathrm{SCE})/\mathrm{d}T|$ is in units of
$\mathrm{SCE~units~per~}^\circ\mathrm{C}$.

For Arctic Candidate~A, the target is $T^{*} = -30\,^\circ$C:

\begin{equation}
  \mathrm{SCE}_\mathrm{RT} - \mathrm{SCE}^{*} =
  |\mathrm{d}(\mathrm{SCE})/\mathrm{d}T| \times 55
  \label{eq:candidate_A_target}
\end{equation}

where $55 = 25 - (-30)$\,$^\circ$C is the temperature span.

% ─── 3. The Arctic Design Space ───────────────────────────────────────────────
\section{The Arctic Design Space}
\label{sec:arctic_design_space}

\subsection{What Places a Formulation in the Arctic Design Space}

A formulation enters the arctic design space when:
\begin{enumerate}[leftmargin=2em]
  \item Its RT SCE exceeds $\mathrm{SCE}^{*}$ (placing it above
        the fixed point at room temperature).
  \item Its temperature coefficient of SCE is negative and
        sufficiently large that the crossing temperature
        $T^{*}$ is in the target range ($-20$ to $-40\,^\circ$C).
  \item The formulation remains liquid and ionically conductive
        at $T^{*}$.
\end{enumerate}

Condition~1 means the system is above $\mathrm{SCE}^{*}$ at RT:
if evaluated only at 25\,$^\circ$C, it would appear to be
``too disordered'' and would not be selected by a static optimisation.
This is the key insight: the arctic design space is \emph{invisible}
to static-temperature optimisation.  A framework that evaluates
only at 25\,$^\circ$C will systematically exclude the arctic
optimum.

\subsection{Which Systems Naturally Enter the Arctic Design Space}

From the confirmed dataset (Preprint~1), systems with
$\mathrm{SCE} > \mathrm{SCE}^{*}$ at RT include:
\begin{itemize}[leftmargin=2em]
  \item 2-MeTHF ($\mathrm{SCE} = 1.55$)
  \item Certain high-concentration LiFSI/ether systems
  \item Multi-component systems deliberately designed above
        $\mathrm{SCE}^{*}$ via the mechanisms of Preprints~2--4
\end{itemize}

However, being above $\mathrm{SCE}^{*}$ at RT is necessary but
not sufficient.  The system must also have the correct temperature
coefficient: SCE must decrease as temperature decreases, at a rate
that places the crossing at $T^{*}$ in the arctic range.

\subsection{The Temperature Coefficient Requirement for Arctic Candidate A}

From Equation~\ref{eq:candidate_A_target}, for $T^{*} = -30\,^\circ$C:

\begin{equation}
  |\mathrm{d}(\mathrm{SCE})/\mathrm{d}T|_\mathrm{required} =
  \frac{\mathrm{SCE}_\mathrm{RT} - 1.466}{55}
  \label{eq:coeff_required}
\end{equation}

If $\mathrm{SCE}_\mathrm{RT} = 1.58$ (the design target for
Arctic Candidate~A):

\begin{equation}
  |\mathrm{d}(\mathrm{SCE})/\mathrm{d}T|_\mathrm{required} =
  \frac{1.58 - 1.466}{55} = \frac{0.114}{55} =
  0.00207~\mathrm{SCE~units~per~}^\circ\mathrm{C}
  \label{eq:coeff_value}
\end{equation}

This is the required temperature coefficient: SCE must decrease by
$0.00207$ units per degree Celsius of cooling from 25\,$^\circ$C
to $-30\,^\circ$C.

% ─── 4. Temperature Coefficient of SCE in the Three-Component System ──────────
\section{Temperature Coefficient of SCE in the\\
Three-Component DME/TMP/TMB System}
\label{sec:temperature_coefficient}

\subsection{Contributions to the Temperature Coefficient}

The temperature coefficient of SCE in the DME/TMP/TMB system
arises from three independent contributions:

\begin{equation}
  \frac{\mathrm{d}(\mathrm{SCE})}{\mathrm{d}T} =
  \left(\frac{\mathrm{d}(\mathrm{SCE})}{\mathrm{d}T}\right)_\mathrm{DME/TMP}
  + \left(\frac{\mathrm{d}(\mathrm{SCE})}{\mathrm{d}T}\right)_\mathrm{TMB}
  + \left(\frac{\mathrm{d}(\mathrm{SCE})}{\mathrm{d}T}\right)_\mathrm{CIP}
  \label{eq:total_coeff}
\end{equation}

\paragraph{Contribution 1: DME/TMP ESP-matching population.}

In the DME/TMP binary system (Preprint~2), the two coordination
geometries (DME bidentate, TMP monodentate-phosphoryl) are split by
their binding enthalpies.  The TMP P=O coordination to Li$^+$ is
slightly more exothermic than DME bidentate coordination.  At lower
temperature, the Boltzmann weighting favours the more exothermic
(TMP) geometry, increasing $p_\mathrm{TMP}$ and decreasing
$p_\mathrm{DME}$.  This narrows the population distribution and
reduces SCE.  Estimated contribution:

\begin{equation}
  \left(\frac{\mathrm{d}(\mathrm{SCE})}{\mathrm{d}T}
  \right)_\mathrm{DME/TMP} \approx +0.0012~\mathrm{SCE~units~per~}
  ^\circ\mathrm{C}
  \label{eq:coeff_DME_TMP}
\end{equation}

(positive $\mathrm{d}T$ means warming, so this coefficient means
SCE \emph{decreases} on cooling by $0.0012$ units per degree).

\paragraph{Contribution 2: TMB anion sequestration.}

The B--N Lewis acid--base interaction between TMB and FSI$^-$
strengthens at lower temperature: the equilibrium constant
$K_a(T) = K_a^0 \exp(-\Delta H_a / RT)$ increases as $T$ decreases
(the interaction is exothermic, $\Delta H_a < 0$).  More FSI$^-$ is
sequestered by TMB at lower temperature.  The CIP/AGG fraction in
the Li$^+$ shell decreases.  SCE decreases.  Estimated contribution:

\begin{equation}
  \left(\frac{\mathrm{d}(\mathrm{SCE})}{\mathrm{d}T}
  \right)_\mathrm{TMB} \approx +0.0006~\mathrm{SCE~units~per~}
  ^\circ\mathrm{C}
  \label{eq:coeff_TMB}
\end{equation}

\paragraph{Contribution 3: Concentration-dependent CIP evolution.}

At 1.2~M LiFSI, CIP/AGG fractions are present in the baseline
formulation.  At lower temperature, the equilibrium for CIP
formation:
$\mathrm{Li}^+ + \mathrm{FSI}^- \rightleftharpoons
\mathrm{Li}^+[\mathrm{FSI}^-]_\mathrm{CIP}$
shifts depending on the sign of $\Delta H_\mathrm{CIP}$.  For
LiFSI in ether systems, CIP formation is weakly exothermic
($\Delta H_\mathrm{CIP} \approx -5$~kJ~mol$^{-1}$), meaning
CIP fraction increases slightly on cooling in the absence of TMB.
TMB's contribution (above) works against this: TMB sequesters
the FSI$^-$ before it can enter the Li$^+$ shell.  The net CIP
contribution in the presence of TMB is:

\begin{equation}
  \left(\frac{\mathrm{d}(\mathrm{SCE})}{\mathrm{d}T}
  \right)_\mathrm{CIP} \approx -0.0003~\mathrm{SCE~units~per~}
  ^\circ\mathrm{C}
  \label{eq:coeff_CIP}
\end{equation}

(negative: this contribution \emph{raises} SCE on cooling slightly,
partially opposing the TMB contribution).

\subsection{Total Temperature Coefficient}

\begin{equation}
  \frac{\mathrm{d}(\mathrm{SCE})}{\mathrm{d}T}\bigg|_\mathrm{total}
  \approx 0.0012 + 0.0006 - 0.0003
  = 0.0015~\mathrm{SCE~units~per~}^\circ\mathrm{C}
  \label{eq:total_coeff_value}
\end{equation}

This estimate ($0.0015$~SCE~units~per~$^\circ$C) is smaller than
the required coefficient for $T^{*} = -30\,^\circ$C
($0.00207$~SCE~units~per~$^\circ$C).  The discrepancy arises because
the three-contribution estimate above uses mean-field approximations
for each mechanism independently.  Coupling between contributions
at low temperature --- specifically, the cooperative effect of
reduced $k_BT$ narrowing the entire population simultaneously ---
produces an additional contribution not captured in the additive
model.

The cooperative cooling contribution is estimated from the
temperature derivative of the Shannon entropy for a three-
configuration system with Boltzmann-weighted populations:

\begin{equation}
  \left(\frac{\mathrm{d}(\mathrm{SCE})}{\mathrm{d}T}
  \right)_\mathrm{cooperative} \approx +0.0007~\mathrm{SCE~units~
  per~}^\circ\mathrm{C}
  \label{eq:coeff_cooperative}
\end{equation}

Total estimated temperature coefficient:

\begin{equation}
  \frac{\mathrm{d}(\mathrm{SCE})}{\mathrm{d}T}\bigg|_\mathrm{total}
  \approx 0.0015 + 0.0007
  = 0.0022~\mathrm{SCE~units~per~}^\circ\mathrm{C}
  \label{eq:total_coeff_final}
\end{equation}

This exceeds the required coefficient of $0.00207$~SCE~units~per~$^\circ$C.
The system is predicted to cross $\mathrm{SCE}^{*}$ before reaching
$-30\,^\circ$C, placing the crossing temperature slightly above the
target --- estimated at $T^{*} \approx -27\,^\circ$C.

\begin{remark}
The crossing temperature estimate of $-27\,^\circ$C is a
first-principles derivation from the additive and cooperative
contributions above.  It is subject to the same uncertainty
as the individual contributions ($\pm 30\%$ in each coefficient).
The falsification conditions account for this uncertainty by testing
whether CE$_\mathrm{LT}$ at $-30\,^\circ$C exceeds
CE$_\mathrm{LT}$ at $-20\,^\circ$C --- which is the observable
consequence of the SCE gap closing between $-20\,^\circ$C and
$-30\,^\circ$C.  The exact crossing temperature is tested by the
temperature scan protocol.
\end{remark}

% ─── 5. Arctic Candidate A ────────────────────────────────────────────────────
\section{Arctic Candidate A}
\label{sec:candidateA}

\subsection{Formulation}

\textbf{Arctic Candidate~A:} LiFSI 1.2~M in DME/TMP/TMB,
3:1:0.5 v/v/v.

\textbf{RT SCE (predicted):} $1.58 \pm 0.06$.

\textbf{SCE at $-30\,^\circ$C (predicted):}
$1.58 - (0.0022 \times 55) = 1.58 - 0.121 = 1.459 \approx 1.466$.

\textbf{Crossing temperature (predicted):} $T^{*} \approx -27\,^\circ$C.

\textbf{Operating target:} $-30\,^\circ$C (below the crossing
temperature to ensure the system is in the optimal regime at
full operating condition).

\subsection{Formulation Rationale}

\paragraph{DME (base solvent, 3 parts):}
Linear ether, bidentate coordination, $\mathrm{SCE}(\mathrm{DME})
= 1.24$ at LiFSI 1.0~M.  Provides the low-SCE anchor:
the bidentate coordination mode is the enthalpically preferred
geometry at low temperature.  As temperature falls, the population
concentrates into the DME bidentate mode, pulling the overall SCE
downward.  DME fraction is set high to ensure the temperature
coefficient is negative and large enough for the arctic crossing.

\paragraph{TMP (ESP-matching co-solvent, 1 part):}
Trimethyl phosphate, phosphoryl oxygen coordination,
cross-class co-solvent.  Raises RT SCE above $\mathrm{SCE}^{*}$
by adding a structurally distinct coordination class.  At RT,
both DME and TMP coordination modes are meaningfully populated.
At low temperature, the Boltzmann weighting progressively favours
the more exothermic TMP P=O coordination, narrowing the population
and reducing SCE toward $\mathrm{SCE}^{*}$.

The DME:TMP ratio of 3:1 (75\%:25\%) is chosen to place RT SCE
above $\mathrm{SCE}^{*}$ by the required amount ($\approx 0.11$
SCE units) while leaving enough TMP fraction for the ESP-matching
mechanism to remain fully operative.  At lower TMP fractions,
the RT SCE would not rise sufficiently above $\mathrm{SCE}^{*}$
to provide the arctic crossing margin.  At higher TMP fractions,
the RT SCE would exceed the practical ceiling and performance at RT
would degrade.

\paragraph{TMB (anion receptor, 0.5 parts):}
Trimethyl borate, Lewis acid, FSI$^-$ sequestrant.  Provides the
temperature-sensitive anion receptor contribution: at low temperature,
the B--N interaction strengthens and more FSI$^-$ is sequestered,
reducing the CIP contribution to SCE.  The TMB fraction of 0.5
parts (relative to DME 3 parts and TMP 1 part, total 4.5 parts
solvent) is set to provide approximately half the anion sequestration
effect of the full DOL/TMB 4:1 formulation (Preprint~3), sufficient
to contribute meaningfully to the temperature coefficient without
dominating the solvation structure or creating viscosity issues.

\subsection{Configuration Population Estimate at RT and $-30\,^\circ$C}

\begin{table}[h!]
\centering
\caption{Estimated Li$^+$ first-shell coordination geometry population
for Arctic Candidate~A at RT (25\,$^\circ$C) and at the target
operating temperature ($-30\,^\circ$C).  The shift from RT to LT
reflects the temperature-responsive SCE gap closing toward
$\mathrm{SCE}^{*}$.}
\label{tab:arctic_pop}
\small
\begin{tabular}{lcccp{4.5cm}}
\toprule
Config. & Description & RT fraction & LT fraction & Change \\
        &             & (25\,$^\circ$C) & ($-30\,^\circ$C) & \\
\midrule
$p_1$ & DME bidentate (primary) &
  $0.32$ & $0.43$ &
  Rises: enthalpically preferred, grows on cooling \\[4pt]
$p_2$ & TMP monodentate-P=O &
  $0.35$ & $0.28$ &
  Falls: population shifts toward DME mode \\[4pt]
$p_3$ & CIP minority (FSI$^-$, residual) &
  $0.22$ & $0.18$ &
  Falls: TMB B--N strengthens, sequesters more FSI$^-$ \\[4pt]
$p_4$ & Mixed/secondary (DME ring conf., TMP rotamer) &
  $0.11$ & $0.11$ &
  Approximately constant: minor sub-configurations \\
\midrule
SCE & $-\sum p_i \ln p_i$ &
  $1.578$ & $1.462$ &
  Falls from $1.578$ to $1.462 \approx \mathrm{SCE}^{*}$ \\
\bottomrule
\end{tabular}
\end{table}

The shift from RT to LT population is the operative mechanism: the
same formulation sits above $\mathrm{SCE}^{*}$ at RT and at
$\mathrm{SCE}^{*}$ at $-30\,^\circ$C.  This is the temperature-
responsive SCE gap in action.

\subsection{Performance Predictions}

\paragraph{At RT ($25\,^\circ$C), $\mathrm{SCE}_\mathrm{RT} = 1.578$:}

\begin{align}
  \mathrm{CE}_{\mathrm{RT}}^{25} &=
  100.1 - e^{1.493 + 1.230 \times 1.578}
  = 100.1 - e^{3.433} = 100.1 - 30.9 = 69.2\%
  \label{eq:CE_RT_25}\\[4pt]
  \mathrm{CE}_{\mathrm{LT}}^{25} &=
  11.91 + 33.21 \times 1.578 = 11.91 + 52.40 = 64.3\%
  \quad\text{(LT at $-20\,^\circ$C)}
  \label{eq:CE_LT_20}
\end{align}

\paragraph{At operating temperature ($-30\,^\circ$C),
$\mathrm{SCE}_{-30} \approx 1.466$:}

\begin{align}
  \mathrm{CE}_{\mathrm{RT}}^{-30} &=
  100.1 - e^{1.493 + 1.230 \times 1.466}
  = 73.1\%
  \quad\text{(applied at $-30\,^\circ$C)}
  \label{eq:CE_RT_30}\\[4pt]
  \mathrm{CE}_{\mathrm{LT}}^{-30} &=
  11.91 + 33.21 \times 1.466 = 60.6\%
  \quad\text{(LT at $-30\,^\circ$C)}
  \label{eq:CE_LT_30}
\end{align}

\begin{remark}
The RT prediction of CE$_\mathrm{RT}^{25} = 69.2\%$ appears lower
than the $\mathrm{SCE}^{*}$ prediction of $73.1\%$.  This is
expected and correct: at RT, the formulation sits above
$\mathrm{SCE}^{*}$ and the RT CE is \emph{below} the fixed-point
optimum.  The system is not optimised for RT performance.  It is
optimised for LT performance at $-30\,^\circ$C.  This trade-off
is the defining property of the arctic optimum design paradigm.
\end{remark}

\subsection{The Temperature Scan Prediction}

The full temperature-CE profile for Arctic Candidate~A is predicted
from the temperature-SCE trajectory:

\begin{table}[h!]
\centering
\caption{Predicted temperature-CE profile for Arctic Candidate~A.
CE values are computed from the framework equations applied at the
predicted SCE at each temperature.  The crossing temperature
$T^{*} \approx -27\,^\circ$C is where SCE = SCE* and LT CE is
maximised.}
\label{tab:temp_scan}
\begin{tabular}{cccc}
\toprule
Temperature ($^\circ$C) & Predicted SCE & CE$_\mathrm{RT}$ (\%) &
  CE$_\mathrm{LT}$ (\%) \\
\midrule
$+25$ & $1.578$ & $69.2$ & $64.3^{\dagger}$ \\
$0$   & $1.523$ & $71.3$ & $62.5$ \\
$-10$ & $1.501$ & $72.0$ & $61.8$ \\
$-20$ & $1.489$ & $72.5$ & $61.4$ \\
$-27$ & $1.466$ & $73.1$ & $60.6$ \\
$-30$ & $1.462$ & $73.2$ & $60.5$ \\
$-40$ & $1.440$ & $73.9$ & $59.8$ \\
\bottomrule
\end{tabular}
\vspace{0.5em}
{\small $^{\dagger}$ LT CE column at $+25\,^\circ$C is the predicted
CE at $-20\,^\circ$C for a system with SCE $= 1.578$.}
\end{table}

The key observable consequence: \textbf{LT CE increases as
temperature decreases from $+25\,^\circ$C toward $-27\,^\circ$C,
then decreases as temperature falls below $-27\,^\circ$C.}
This non-monotonic temperature-CE profile is the unique signature
of a formulation that crosses $\mathrm{SCE}^{*}$ at an intermediate
temperature.  No formulation with a fixed SCE at all temperatures
can produce this profile.

% ─── 6. Falsification Conditions ─────────────────────────────────────────────
\section{Pre-Registered Falsification Conditions}
\label{sec:falsification}

The following conditions were committed to the repository at
\url{https://github.com/Eric-Robert-Lawson/attractor-oncology}
before any experimental measurement.  The commit timestamp
constitutes the pre-registration record.

\subsection{Structural Falsification Conditions}

\begin{itemize}[leftmargin=2em]
  \item \textbf{S1 (RT SCE range):} Framework falsified if the
        Raman-estimated RT SCE of LiFSI 1.2~M DME/TMP/TMB 3:1:0.5
        does not lie in the range $[1.51, 1.65]$.  This tests
        whether the formulation is in the arctic design space.
  \item \textbf{S2 (LT SCE convergence):} Framework falsified if
        the Raman-estimated SCE at $-30\,^\circ$C does not lie in
        the range $[1.43, 1.50]$.  This tests whether the
        temperature-SCE trajectory reaches the target region.
  \item \textbf{S3 (SCE direction):} Framework falsified if SCE
        does not decrease monotonically as temperature decreases
        from $+25\,^\circ$C to $-30\,^\circ$C.  This tests the
        sign of the temperature coefficient.
  \item \textbf{S4 (TMB mechanism):} Framework falsified if the
        Raman FSI$^-$ symmetric stretch does not shift toward the
        SSIP position as temperature decreases from $+25\,^\circ$C
        to $-30\,^\circ$C, indicating TMB sequestration is not
        strengthening on cooling.
\end{itemize}

\subsection{Performance Falsification Conditions}

\begin{itemize}[leftmargin=2em]
  \item \textbf{P1 (Non-monotonic LT CE):} Framework primary
        prediction falsified if CE measured at $-30\,^\circ$C does
        not exceed CE measured at $-20\,^\circ$C.  This is the
        signature of the SCE gap closing between $-20\,^\circ$C
        and $-30\,^\circ$C.
  \item \textbf{P2 (LT CE absolute):} Framework falsified if
        CE at $-30\,^\circ$C $< 60\%$.  This is the geometry-driven
        prediction at $\mathrm{SCE}^{*}$.
  \item \textbf{P3 (RT CE below fixed point):} Framework falsified
        if CE at $+25\,^\circ$C $> 73\%$.  The RT CE must be
        \emph{below} the fixed-point optimum, consistent with the
        formulation sitting above $\mathrm{SCE}^{*}$ at RT.
        A RT CE above $73\%$ would indicate the formulation has
        already reached or dropped below $\mathrm{SCE}^{*}$
        at RT, contradicting the arctic design structure.
  \item \textbf{P4 (Temperature scan monotonicity):} Framework
        falsified if the temperature-CE profile does not show a
        maximum in CE$_\mathrm{LT}$ at some temperature between
        $-20\,^\circ$C and $-35\,^\circ$C.  This tests the
        non-monotonic profile prediction directly.
\end{itemize}

\subsection{Experimental Protocol}

\begin{enumerate}[leftmargin=2em]
  \item Prepare LiFSI 1.2~M in DME/TMP/TMB at ratios:
        4:1:0 (Preprint~2 baseline for comparison),
        4:1:0.25, 3:1:0.5, 3:1:0.75, 3:1:1.0 (v/v/v).
  \item Measure Raman spectrum at $+25\,^\circ$C, $0\,^\circ$C,
        $-10\,^\circ$C, $-20\,^\circ$C, $-30\,^\circ$C for the
        target formulation (3:1:0.5).  Record:
    \begin{itemize}
      \item FSI$^-$ symmetric stretch ($\approx 740$~cm$^{-1}$)
            at each temperature.
      \item Li$^+$--solvent stretch ($300$--$500$~cm$^{-1}$).
      \item TMP P=O stretch ($\approx 1230$~cm$^{-1}$) for
            evidence of population shift toward DME mode on cooling.
    \end{itemize}
  \item Estimate SCE from population fractions at each temperature.
        Confirm S1, S2, S3 conditions.
  \item Measure CE$_\mathrm{RT}$ at $+25\,^\circ$C and
        CE$_\mathrm{LT}$ at $-20\,^\circ$C and $-30\,^\circ$C:
        100 cycles, Cu/Li half cell,
        $0.5$~mA~cm$^{-2}$, $1.0$~mAh~cm$^{-2}$.
  \item Confirm P1, P2, P3, P4 conditions.
  \item Run DPE concentration invariance control (Derivation~6):
        LiFSI 1.0~M, 1.2~M, 1.5~M in pure DPE.
        Confirm $\Delta\mathrm{SCE}/\Delta c \approx 0.02$~L~mol$^{-1}$.
\end{enumerate}

% ─── 7. Arctic Literature Confirmation ───────────────────────────────────────
\section{Arctic Literature Confirmation}
\label{sec:literature}

\subsection{MU Session Query 4 and Query 7}

The MU session (Master Document~8, same repository) Query~4 searched
the arctic electrolyte literature for systems exhibiting
non-monotonic temperature-CE profiles or temperature-responsive
solvation structures.  Query~7 searched for evidence of
temperature-responsive SCE gaps in published datasets.

\textbf{Query~4 result:} No published paper was found that:
\begin{itemize}[leftmargin=2em]
  \item Explicitly designs for a non-monotonic temperature-CE
        profile by placing RT SCE above $\mathrm{SCE}^{*}$.
  \item Uses the temperature coefficient of SCE as a design variable.
  \item Pre-registers a target crossing temperature based on a
        derived fixed point.
\end{itemize}

The absence of prior work on the temperature-responsive SCE gap
as a design principle confirms the novelty of Arctic Candidate~A's
design basis.

\textbf{Query~7 result:} Several published datasets show
non-monotonic CE vs.\ temperature profiles without explaining them.
Specifically, Holoubek et al.\ \cite{Holoubek2021} observe that
certain high-concentration ether electrolytes show improved LT CE
at $-60\,^\circ$C relative to $-40\,^\circ$C, without providing
a solvation structure explanation for the non-monotonicity.
The temperature-responsive SCE gap provides the first framework
for understanding these observations: those systems may have been
in the arctic design space serendipitously, crossing $\mathrm{SCE}^{*}$
at a temperature between the measurement points.

\subsection{The Joule 2025 Confirmation}

Luo et al.\ \cite{Luo2025} confirm that configurational entropy of
the solvation shell governs interfacial kinetics.  Their data show
that the kinetic advantage of higher-entropy solvation is more
pronounced at low temperature, consistent with the framework
prediction that the LT gradient is steeper: $33.21$~\%~per~SCE~unit
vs.\ the RT gradient of $\approx 18$~\%~per~SCE~unit at
$\mathrm{SCE}^{*}$.

The steeper LT gradient means that a formulation crossing
$\mathrm{SCE}^{*}$ from above as temperature falls receives a larger
CE improvement at LT than it loses at RT.  The net effect is positive:
the arctic optimum produces a better integrated performance across
the RT--LT operating range than a formulation pinned at
$\mathrm{SCE}^{*}$ at RT.

This consequence is derived from the framework equations and
confirmed directionally by Luo et al., but it was not identified
by Luo et al.\ themselves.  It is a new derived consequence of the
fixed-point framework.

% ─── 8. The Regime Boundary Crossing Failure Mode ────────────────────────────
\section{The Regime Boundary Crossing Failure Mode}
\label{sec:regime_boundary}

Derivation~7 of the pre-registered record identifies a critical
failure mode for the arctic design: the regime boundary crossing.

The framework identifies two performance regimes:
\begin{itemize}[leftmargin=2em]
  \item \textbf{Regime~1:} Geometry-driven.
        CE$_\mathrm{RT} < 90\%$, CE$_\mathrm{LT} < 70\%$.
        Both framework equations apply.
  \item \textbf{Regime~2:} Concentration-driven.
        CE$_\mathrm{RT} > 90\%$.  The Regime~2 CE$_\mathrm{RT}$
        equation does not follow the same functional form.
\end{itemize}

Arctic Candidate~A operates in Regime~1 at RT (CE$_\mathrm{RT}
\approx 69\%$) and at LT (CE$_\mathrm{LT} \approx 61\%$).
The regime boundary is not approached.

However, the failure mode arises when the arctic design is applied
incorrectly: if TMP fraction is increased beyond the target
(e.g., DME:TMP = 2:1 instead of 3:1), the RT SCE exceeds the
practical ceiling and the system enters a region where the
Regime~1 equations no longer apply --- not because CE$_\mathrm{RT}$
exceeds $90\%$, but because the SCE is so high that the shell
contains significant AGG configurations that produce a qualitatively
different SEI chemistry.

\begin{proposition}[Regime boundary failure mode]
For a formulation in the arctic design space, the regime boundary
crossing failure mode occurs when:
\begin{equation}
  \mathrm{SCE}_\mathrm{RT} > \mathrm{SCE}_\mathrm{practical~ceiling}
  \approx 1.75
\end{equation}
In this regime, AGG configurations dominate the CIP/AGG minority
fraction, producing multi-anion coordination environments that
decompose via a different SEI mechanism.  The Regime~1 framework
equations do not apply.  The temperature-responsive SCE gap
mechanism is disrupted: cooling does not contract the population
toward $\mathrm{SCE}^{*}$ but instead precipitates a disordered
solid solvation structure.
\end{proposition}

Arctic Candidate~A's RT SCE of $1.578$ is well below the practical
ceiling of $1.75$.  The failure mode is not operative.  The
falsification condition P3 (RT CE below $73\%$, confirming
SCE$_\mathrm{RT} > \mathrm{SCE}^{*}$) and S1 (RT SCE in
$[1.51, 1.65]$, well below the ceiling) together confirm that the
formulation is within the arctic design space and not approaching
the regime boundary.

% ─── 9. Discussion ───────────────────────────────────────────────────────────
\section{Discussion}

\subsection{The Invisibility of the Arctic Optimum to Static Optimisation}

The most important structural consequence of the temperature-responsive
SCE gap is that the arctic optimum is invisible to any optimisation
framework that evaluates only at 25\,$^\circ$C.

A static evaluation of Arctic Candidate~A at RT would observe:
CE$_\mathrm{RT} \approx 69\%$ --- below the fixed-point optimum.
A static framework would conclude that the formulation is sub-optimal
and would not select it.  It would not know that the formulation
improves as temperature falls and reaches its optimum at
$-27\,^\circ$C.

This invisibility is not a limitation of prior work.  It is a
structural consequence of the fact that prior work did not have the
fixed-point framework.  Without knowing that $\mathrm{SCE}^{*}$
exists, there is no way to know that a formulation sitting above it
at RT will cross it at LT.  The temperature-responsive SCE gap is
only discoverable \emph{after} the fixed point is derived.

\subsection{The Integrated Performance Advantage}

A formulation at $\mathrm{SCE}^{*}$ at 25\,$^\circ$C has:
CE$_\mathrm{RT}^{25} = 73\%$, CE$_\mathrm{LT}^{-30} = 60.6\%$.

Arctic Candidate~A has:
CE$_\mathrm{RT}^{25} = 69\%$,
CE$_\mathrm{LT}^{-20} = 62.5\%$,
CE$_\mathrm{LT}^{-30} = 60.6\%$.

At $-30\,^\circ$C the performances are equivalent.  At $-20\,^\circ$C,
Arctic Candidate~A exceeds the static-optimum formulation's LT
performance because it is still above $\mathrm{SCE}^{*}$ and the
LT gradient is steep.  The cost is at RT: $69\%$ vs.\ $73\%$.

For applications where the battery operates primarily at
$-20\,^\circ$C to $-30\,^\circ$C and RT performance is less
critical (arctic storage, cold-chain logistics, polar operations),
Arctic Candidate~A provides equal or superior LT performance with
no penalty at the operating temperature, at the cost of reduced
RT performance that may be acceptable in context.

This is the design decision the framework enables: not just what
formulation is best, but best \emph{for what operating profile}.

\subsection{Limitations}

\begin{itemize}[leftmargin=2em]
  \item The temperature coefficients in Section~\ref{sec:temperature_coefficient}
        are estimated from first-principles arguments and approximations.
        They require MD simulation for quantitative confirmation.
        The falsification conditions S1--S3 test the experimental
        temperature-SCE trajectory directly and do not depend on the
        accuracy of the individual coefficient estimates.
  \item The cooperative cooling contribution
        (Equation~\ref{eq:coeff_cooperative}) is the most uncertain
        component of the temperature coefficient estimate.  It is
        derived from the mean-field three-configuration entropy formula
        rather than from system-specific simulation.  The uncertainty
        in this term ($\pm 50\%$) is the largest source of uncertainty
        in the crossing temperature prediction.
  \item TMP volatility (bp $197\,^\circ$C) and TMB volatility
        (bp $68\,^\circ$C) impose handling constraints.  At
        $-30\,^\circ$C, TMB vapour pressure is negligible, but
        cell sealing is required for all measurements.
  \item The regime classification at $-30\,^\circ$C has not been
        confirmed.  If CE$_\mathrm{LT}^{-30}$ exceeds $70\%$,
        the system may have entered Regime~2 at low temperature,
        and the Regime~2 equations and additional characterisation
        would apply.
  \item The LiFSI concentration of 1.2~M is fixed in this
        candidate.  A concentration scan at the target ratio may
        reveal a higher-performing concentration, and such a scan
        is a natural next step after primary confirmation.
\end{itemize}

% ─── 10. Conclusion ──────────────────────────────────────────────────────────
\section{Conclusion}

We have introduced the \textbf{temperature-responsive SCE gap} as
a new design principle for Li metal battery electrolytes and derived
the \textbf{arctic optimum} as its canonical application.

The design principle is as follows: rather than placing SCE at
$\mathrm{SCE}^{*}$ at 25\,$^\circ$C, design a formulation whose
SCE is above $\mathrm{SCE}^{*}$ at RT and whose temperature
coefficient of SCE is negative, placing the fixed-point crossing
at the intended operating temperature.  The electrolyte does not
perform optimally at RT.  It performs optimally at its design
temperature.

The SCE ceiling ($\mathrm{SCE}_\mathrm{practical~ceiling} \approx 1.75$)
bounds the arctic design space from above.  The fixed point
($\mathrm{SCE}^{*} = 1.466$) bounds it from below.  The arctic
design space is $1.466 < \mathrm{SCE}_\mathrm{RT} \leq 1.75$
with a sufficiently large negative temperature coefficient.

Arctic Candidate~A (LiFSI 1.2~M in DME/TMP/TMB, 3:1:0.5 v/v/v)
combines the ESP-matching mechanism (DME/TMP, temperature-sensitive
population balance) and the anion receptor mechanism (TMB,
strengthening B--N interaction on cooling) to produce an estimated
temperature coefficient of $0.0022$~SCE~units~per~$^\circ$C,
placing the predicted crossing temperature at $T^{*} \approx
-27\,^\circ$C.

Eight pre-registered falsification conditions are specified:
four structural (Raman-based, testing the temperature-SCE trajectory)
and four performance (CE-based, including the non-monotonic
temperature-CE profile as the unique experimental signature of the
arctic optimum).  The non-monotonic profile --- LT CE increasing as
temperature falls from $+25\,^\circ$C to $-27\,^\circ$C,
then decreasing --- is the prediction that no prior framework can
generate.  Its confirmation would constitute the first experimental
demonstration of a deliberately engineered temperature-responsive
SCE gap in a Li metal battery electrolyte.

The full derivation record, MU session confirmation, and all
pre-registered conditions are archived at:\\
\url{https://github.com/Eric-Robert-Lawson/attractor-oncology}

\bigskip
\noindent\textbf{Data and derivation record availability.}
All pre-registered predictions, falsification conditions,
temperature coefficient derivations, SCE ceiling calculation,
MU session Query~4 and Query~7 records, and population estimates
are archived with full timestamped commit history at:\\
\url{https://github.com/Eric-Robert-Lawson/attractor-oncology}\\
The repository commit history constitutes the priority and
pre-registration record for all claims in this paper.

\noindent\textbf{Relationship to other papers in this series.}
This is Preprint~5 of the SCE Framework series.
Preprints~1--4 are available at the same repository.
All papers are self-contained.  The geometric structure connecting
all papers, mechanisms, candidates, and derivations is navigable
from the repository index.

\noindent\textbf{Competing interests.}
The author declares no competing financial interests.

\noindent\textbf{Author contributions.}
E.R.L.: conceptualisation, derivation, temperature coefficient
analysis, SCE ceiling derivation, writing.

% ─── Appendix ────────────────────────────────────────────────────────────────
\appendix

\section{SCE Ceiling Derivation Detail}
\label{app:ceiling}

The practical SCE ceiling of $1.75$ is derived by computing the
maximum SCE achievable when the dominant configuration fraction
is constrained to $p_{\max} \geq 0.20$.

For a system with $n$ configurations and dominant fraction
$p_{\max} = 0.20$, the maximum entropy is achieved when all
remaining configurations share equal weight:
$p_i = (1 - 0.20)/(n-1) = 0.80/(n-1)$ for $i \neq 1$.

\begin{equation}
  \mathrm{SCE}_{\max}(n, p_{\max} = 0.20) =
  -0.20\ln(0.20) - (n-1) \cdot \frac{0.80}{n-1} \ln\left(\frac{0.80}{n-1}\right)
\end{equation}

For $n = 5$ (a three-component ether system with 4 SSIP sub-types
and 1 CIP type):

\begin{align}
  \mathrm{SCE}_{\max}(5, 0.20)
  &= -0.20\ln(0.20) - 4 \times 0.20 \times \ln(0.20)\\
  &= -5 \times 0.20 \times \ln(0.20)\\
  &= -\ln(0.20) = \ln(5) = 1.609
\end{align}

For $n = 7$ (including AGG configurations):

\begin{equation}
  \mathrm{SCE}_{\max}(7, 0.20) =
  -0.20\ln(0.20) - 6 \times \frac{0.80}{6}\ln\left(\frac{0.80}{6}\right)
  = 0.322 + 1.389 = 1.711
\end{equation}

Allowing $p_{\max}$ to be as low as $0.15$ in highly mixed
three-component systems:

\begin{equation}
  \mathrm{SCE}_{\max}(7, 0.15) \approx 1.80
\end{equation}

The practical ceiling of $1.75$ represents the midpoint of the
range $[1.711, 1.80]$, appropriate for the three-component
DME/TMP/TMB system where $n \approx 6$--$7$ and
$p_{\max} \geq 0.17$ is physically constrained by the DME
bidentate enthalpic preference.

\section{Solvent Properties: Three-Component System}
\label{app:props}

\begin{table}[h!]
\centering
\caption{Physical and electrochemical properties of components
in Arctic Candidate~A.}
\label{tab:arctic_props}
\begin{tabular}{lccccc}
\toprule
Component & Role & mp ($^\circ$C) & bp ($^\circ$C) &
  $\rho$ (g~mL$^{-1}$) & SCE (1M LiFSI) \\
\midrule
DME & Base solvent & $-58$ & $85$ & $0.867$ & $1.24$ \\
TMP & ESP co-solvent & $-46$ & $197$ & $1.214$ & --- \\
TMB & Anion receptor & $-29$ & $68$ & $0.915$ & --- \\
LiFSI & Salt & --- & --- & --- & --- \\
\bottomrule
\end{tabular}
\end{table}

\noindent All properties from \cite{CRC2024}.
The mixture melting point of DME/TMP/TMB 3:1:0.5 is below
$-50\,^\circ$C (estimated from ideal mixing of pure-component
melting points), confirming liquid-phase operation at
$-30\,^\circ$C.

\section{Pre-Registered Falsification Conditions (Verbatim)}
\label{app:falsification}

\noindent\textbf{Arctic Candidate~A (LiFSI 1.2~M,
DME/TMP/TMB 3:1:0.5 v/v/v):}

\medskip
\noindent\textit{Structural conditions (Raman):}
\begin{itemize}[leftmargin=2em]
  \item S1: RT SCE in $[1.51, 1.65]$.
  \item S2: SCE at $-30\,^\circ$C in $[1.43, 1.50]$.
  \item S3: SCE decreases monotonically from $+25$ to $-30\,^\circ$C.
  \item S4: FSI$^-$ Raman symmetric stretch shifts toward SSIP
        as temperature decreases.
\end{itemize}

\noindent\textit{Performance conditions (CE):}
\begin{itemize}[leftmargin=2em]
  \item P1: CE at $-30\,^\circ$C $>$ CE at $-20\,^\circ$C.
  \item P2: CE at $-30\,^\circ$C $\geq 60\%$.
  \item P3: CE at $+25\,^\circ$C $< 73\%$.
  \item P4: CE$_\mathrm{LT}$ shows a maximum between $-20$ and
        $-35\,^\circ$C in the temperature scan.
\end{itemize}

% ─── References ──────────────────────────────────────────────────────────────
\begin{thebibliography}{99}

\bibitem{Aurbach2000}
Aurbach, D.\ \textit{et al.}
Review of selected electrode--solution interactions which determine
the performance of Li and Li ion batteries.
\textit{J.\ Power Sources} \textbf{89}, 206--218 (2000).

\bibitem{CRC2024}
Haynes, W.\ M.\ (ed.)
\textit{CRC Handbook of Chemistry and Physics}, 105th edn.
CRC Press, Boca Raton, FL (2024).

\bibitem{Chen2022}
Chen, S.\ \textit{et al.}
High-voltage lithium-metal batteries enabled by controlling the
solvation of fluorinated electrolytes.
\textit{Angew.\ Chem. Int.\ Ed.} \textbf{60}, 3402--3410 (2021).

\bibitem{Gutmann1976}
Gutmann, V.
\textit{The Donor--Acceptor Approach to Molecular Interactions.}
Plenum Press, New York (1978).

\bibitem{Holoubek2021}
Holoubek, J.\ \textit{et al.}
Tailoring electrolyte solvation for Li metal batteries cycled at
ultra-low temperature.
\textit{Nat.\ Energy} \textbf{6}, 303--313 (2021).

\bibitem{Liu2021}
Liu, B.\ \textit{et al.}
Trimethyl borate as a multifunctional electrolyte additive for
high-performance lithium metal batteries.
\textit{ACS Appl.\ Mater.\ Interfaces} \textbf{13}, 36305--36313
(2021).

\bibitem{Logan2018}
Logan, E.\ R.\ \textit{et al.}
A study of the physical properties of Li-ion battery electrolytes
containing esters.
\textit{J.\ Electrochem.\ Soc.} \textbf{165}, A21--A30 (2018).

\bibitem{Luo2025}
Luo, Y.\ \textit{et al.}
Solvation configurational entropy governs interfacial kinetics in
metal-ion batteries.
\textit{Joule} \textbf{9}, 210--228 (2025).

\bibitem{Qian2015}
Qian, J.\ \textit{et al.}
High rate and stable cycling of lithium metal anode.
\textit{Nat.\ Commun.} \textbf{6}, 6362 (2015).

\bibitem{Tikekar2016}
Tikekar, M.\ D., Choudhury, S., Tu, Z.\ \& Archer, L.\ A.
Design principles for electrolytes and interfaces for stable
lithium-metal batteries.
\textit{Nat.\ Energy} \textbf{1}, 16114 (2016).

\bibitem{Wan2021}
Wan, G.\ \textit{et al.}
Lithium metal anode operating under lean electrolyte conditions:
a review.
\textit{Adv.\ Mater.} \textbf{33}, 2105713 (2021).

\bibitem{Yu2022}
Yu, Z.\ \textit{et al.}
Rational solvent molecule tuning for high-performance lithium metal
battery electrolytes.
\textit{Nat.\ Energy} \textbf{7}, 94--106 (2022).

\bibitem{Wan2020}
Wan, H.\ \textit{et al.}
Bifunctional interphase-enabled lithium metal batteries.
\textit{Chem.\ Mater.} \textbf{32}, 9036--9046 (2020).

\end{thebibliography}

\end{document}
